% from ega4-2-fr.tex
\oldpage[IV-2]{16}
(i) Si $y\in W$, les points génériques $\eta_i$ ($1\leq i\leq m$) des
composantes irréductibles $Y_i$ de $Y$ contenant $y$ appartiennent par
hypothèse à $W$, donc, puisque $W$ est fermé, ces composantes elles-mêmes sont
contenues dans $W$~; comme il y a par hypothèse un voisinage $U$ de $y$ tel que
$U$ soit réunion des $U\cap Y_i$ $(0_{\mathrm I},2.1.6)$, on a
$U\subset W$, donc $W$ est ouvert. Comme pour tout $y\in Y$,
$\overline{\{y\}}$ est connexe, une composante connexe de $Y$ est stable par
générisation, donc la seconde assertion résulte aussitôt de la première.

(ii) Comme $C$ est fermé dans $X$, $f(C)$ est fermé dans $Y$ par hypothèse. En
outre, comme $C$ est stable par générisation, l'hypothèse sur $f$ entraîne que
$f(C)$ est stable par générisation~; donc, en vertu de (i), $f(C)$ est ouvert
et fermé, et comme il est connexe c'est une composante connexe de $Y$.

(iii) Soit $x\in X$~; par hypothèse, il y a un voisinage ouvert $U$ de
$y=f(x)$ ne rencontrant qu'un nombre fini de composantes irréductibles $Y_i$
de $Y$ ($1\leq i\leq n$)~; soit $y_i$ le point générique de $Y_i$
($1\leq i\leq n$). Pour toute composante irréductible $Z$ de $X$ rencontrant
$f^{-1}(U)$, le point générique $z$ de $Z$ est nécessairement dans l'un des
ensembles $f^{-1}(y_i)$ (2.3.4). Comme chacun de ces ensembles est fini par
hypothèse, cela prouve notre assertion.

\begin{proposition}[2.3.7]
\label{IV.2.3.7-fr}
Soient $f:X\to Y$ un morphisme, $g:Y'\to Y$ un morphisme quasi-plat,
$X'=X\times_Y Y'$, $f'=f_{(Y')}:X'\to Y'$.
\begin{enumerate}
  \item[(i)] Si $f$ est quasi-compact et dominant, il en est de même de $f'$.
  \item[(ii)] Si toute composante irréductible de $X$ domine une composante
  irréductible de $Y$, alors toute composante irréductible de $X'$ domine une
  composante irréductible de $Y'$.
\end{enumerate}
\end{proposition}

Désignons par $g':X'\to X$ la projection canonique, qui est un morphisme
quasi-plat (2.3.3).

(i) On sait déjà (1.1.2) que $f'$ est quasi-compact~; en outre, si $y'$ est un
point maximal (1.1.4) de $Y'$, $y=g(y')$ est point maximal de $Y$, comme il
résulte de (2.3.4, (iii)). Par hypothèse (1.1.5), il existe $x\in X$ tel que
$f(x)=y$, donc (I, 3.4.7), il existe $x'\in X'$ tel que $f'(x')=y'$.

(ii) Soit $x'$ un point maximal de $X'$, et soit $x=g'(x')$~; il résulte de
(2.3.4, (ii)) que $x$ est point maximal de $X$, et par hypothèse $y=f(x)$ est
point maximal de $Y$. Posons $y'=f'(x')$, de sorte que $g(y')=y$~; il s'agit
de montrer que $y'$ est point maximal de $Y'$. En vertu de (I, 5.1.7) et
(2.3.3), on peut se borner au cas où $X$ et $Y$ sont \emph{réduits}, donc
$\mathcal O_x$ et $\mathcal O_y$ sont des \emph{corps}~; en outre, en vertu de
(I, 3.6.5) appliqué deux fois et de (I, 2.4.4), on peut se borner au cas où
$Y=\operatorname{Spec}(\mathcal O_y)$ et
$Y'=\operatorname{Spec}(\mathcal O_{y'})$. Alors $f$ est un morphisme
\emph{plat} puisque $\mathcal O_y$ est un corps (2.1.2), et il en est de même
de $f'$ (2.1.4)~; donc il résulte de (2.3.4, (ii)) que $y'=f'(x')$ est point
maximal de $Y'$.

\begin{corollary}[2.3.8]
\label{IV.2.3.8-fr}
\emph{(Zariski).} Soient $A$, $B$ deux anneaux locaux noethériens,
$\mathfrak m$, $\mathfrak n$ leurs idéaux maximaux respectifs,
$\varphi:A\to B$ un homomorphisme local. On suppose vérifiées les hypothèses
suivantes~:
\begin{enumerate}
  \item[$1^\circ$] $B$ est une $A$-algèbre essentiellement de type fini
  (1.3.8).
  \item[$2^\circ$] Le complété $\widehat A$ de $A$ pour la topologie
  $\mathfrak m$-adique est intègre.
  \item[$3^\circ$] $\varphi$ est injectif.
\end{enumerate}
Alors la topologie $\mathfrak m$-adique de $A$ est induite par la topologie
$\mathfrak n$-adique de $B$.
\end{corollary}

